\section{Evaluación}
\label{sec:evaluacion}

La evaluación pregunta si el motor hace lo que se diseñó para hacer. No pregunta si adivina una compra que nadie observó. Todas las cifras de esta sección salen del archivo del \corte, salvo cuando el texto dice que un resultado es histórico.

Hay cuatro comprobaciones. La primera revisa la regla del universo puntuable. La segunda contrasta trece productos revisados a mano. La tercera pregunta si la parte nutricional se mueve en un sentido razonable frente a una medida externa. La cuarta observa si el orden cambia cuando cambian las prioridades. Después se comprueba que el resultado no depende del azar, y se dice con claridad qué no queda demostrado.

\subsection{¿Se cumple la regla del universo puntuable?}

La regla se cumple sin diferencias: 5\,864 productos tienen procesamiento y al menos cuatro de ocho percentiles. La tabla~\ref{tab:cobertura} separa otros conteos que a veces se confunden con esa regla. Tener la señal nutricional, por ejemplo, no es lo mismo que pertenecer al universo puntuable.

\begin{table}[ht]
  \centering
  \caption{Cobertura del catálogo, antes de aplicar un perfil. Corte del \corte.}
  \label{tab:cobertura}
  \begin{tabular}{lrr}
    \toprule
    Condición & Productos & Porcentaje \\
    \midrule
    Categoría de referencia & 7\,568 & 57{,}80 \\
    Al menos un percentil del núcleo & 7\,047 & 53{,}82 \\
    Dimensión nutricional calculable & 7\,028 & 53{,}68 \\
    Al menos cuatro percentiles & 6\,998 & 53{,}45 \\
    Los ocho percentiles & 6\,083 & 46{,}46 \\
    Procesamiento disponible & 6\,780 & 51{,}78 \\
    Universo puntuable & 5\,864 & 44{,}79 \\
    \bottomrule
  \end{tabular}
\end{table}

\subsection{¿Los casos revisados producen el comportamiento esperado?}

Trece productos reales del catálogo se revisaron a mano, con el resultado esperado de cada caso. Sobre el archivo de este trabajo, doce coinciden. El caso \texttt{5060323907641} difiere solo en el nombre: la revisión esperaba el campo original, y el catálogo ya muestra el nombre tomado del nombre genérico, con el campo original vacío. El puntaje de ese producto sí se calcula. No falla ninguna de las tres señales. El caso no se quitó ni se reescribió para que la cuenta quedara en trece de trece.

\subsection{¿La dimensión nutricional se relaciona con un contraste externo?}

Nutri-Score es una etiqueta nutricional usada como contraste externo. No es la respuesta correcta del motor ni la meta del orden \citep{julia2017nutriscore}. En 6\,035 productos que tienen las dos medidas, la correlación de Spearman entre la señal nutricional y \texttt{nutriscore\_score} es $-0{,}475683$. Spearman resume si dos ordenamientos van en el mismo sentido o en sentidos opuestos.

El signo negativo es el que corresponde. En esta comparación, un valor alto indica mayor compatibilidad con las prioridades nutricionales declaradas. En Nutri-Score, la letra A es el extremo más favorable de esa escala y la E el menos favorable; el puntaje numérico, en cambio, es peor cuando sube. Una relación negativa es coherente con esas dos direcciones. No convierte a Nutri-Score en la verdad del sistema. El proyecto pidió, como mínimo de sanidad, una correlación menor o igual a $-0{,}3$. La cifra medida supera ese umbral en magnitud.

El promedio de la señal nutricional baja de la letra A a la letra E: 59{,}68, 54{,}80, 50{,}87, 49{,}00 y 41{,}57 (figura~\ref{fig:nutriscore}). Esa bajada comprueba la misma dirección.

\begin{figure}[ht]
  \centering
  \includegraphics[width=0.86\textwidth]{d1_nutriscore.png}
  \caption{Dimensión nutricional promedio según la letra de Nutri-Score. \(n = 6\,035\). El título interno de la gráfica llama D1 a esa dimensión. Fuente: notebook maestro, corte del \corte.}
  \label{fig:nutriscore}
\end{figure}

El mismo contraste, calculado el 19 de septiembre de 2026 sobre otro archivo, dio \(n = 6\,041\) y una correlación de $-0{,}4747$. Es un resultado histórico. No sustituye la cifra de este corte.

\subsection{¿El resultado cambia cuando cambian las prioridades?}

El perfil de ejemplo evita gluten, declara dieta vegana, valora las etiquetas orgánico y sin gluten, y pone la nutrición en primer lugar, las etiquetas en segundo y el procesamiento en tercero. No es una persona observada. Con la regla de bandas de la aplicación, el catálogo queda así: 3\,808 excluidos, 9\,171 no verificables, 1 con información insuficiente y 113 en la banda de ranking. Si el procesamiento pasa al primer lugar y la nutrición al tercero, el orden tiene 114 productos. Los excluidos y los no verificables no cambian, porque no dependen de los pesos.

\begin{figure}[ht]
  \centering
  \includegraphics[width=0.86\textwidth]{scores_perfil.png}
  \caption{Resultados del orden cuando la nutrición es la primera prioridad. \(n = 113\). Fuente: notebook maestro, corte del \corte.}
  \label{fig:scores}
\end{figure}

Dentro de los 113, el primer producto obtiene 93{,}30. El cuarto, con resultado 83{,}33, no pertenece al universo puntuable: no tiene señal nutricional y entra porque el procesamiento y las etiquetas cubren la mitad del peso. Al subir el peso del procesamiento, ese producto sale de los diez primeros. El orden responde al perfil.

El cuaderno también mide una regla distinta, solo para ver qué tan sensible es el conteo. En esa regla experimental el producto entra si la alergia es apta, la dieta es compatible o no verificable, y el resultado no es nulo. Con ella hay 665 productos si la nutrición va primero y 640 si va primero el procesamiento. Esas cifras no son la banda de ranking de la aplicación. Confundir 665 y 640 con 113 y 114 mezcla dos definiciones.

\subsection{Determinismo}

Ejecutar dos veces el mismo archivo, el mismo perfil y el mismo conjunto de productos devuelve el mismo resultado, la misma banda y el mismo orden de los 113 productos. El motor no usa azar.

\subsection{Lo que la evaluación demuestra y lo que no}

Con este archivo y este perfil, la evaluación sostiene una afirmación acotada. La regla del universo puntuable se cumple. Doce de trece casos revisados se comportan como se esperaba. La señal nutricional guarda la dirección del contraste con Nutri-Score. El orden cambia cuando cambian las prioridades. La corrida es determinista.

No hay un porcentaje de aciertos contra compras, porque esas compras no existen. No hay un conjunto de prueba separado de un entrenamiento, porque no hay entrenamiento. Nutri-Score permanece como contraste externo, no como la respuesta correcta. La evaluación no dice que la persona vaya a comprar el producto que quedó primero.

Lo que sí queda es un orden que se puede explicar. La sección siguiente muestra cómo ese orden llega a la compra.
